\documentclass[11pt,a4paper]{article}
\usepackage[margin=20mm]{geometry}
\usepackage[T1]{fontenc}
\usepackage[utf8]{inputenc}
\usepackage{graphicx}
\usepackage{hyperref}
\usepackage{float}
\setlength{\parindent}{0pt}
\begin{document}
\begin{titlepage}
\centering
{\LARGE Unified vehicle calibration report\par}
\vspace{1cm}
{\large Session: \texttt{real\_steering\_20260625}\par}
\vfill
This document is regenerated after every completed stage. Each following page records the capture, analysis/update, and independent-validation role of one stage.
\end{titlepage}
\tableofcontents
\clearpage
\clearpage
\small
\section*{Steering command-chain audit}
\addcontentsline{toc}{section}{Steering command-chain audit}
\textbf{Stage key:} \texttt{steering\_command\_audit}\\
\textbf{Status:} pending\\

\subsection*{A --- Capture}
A — Record the stage-specific data.

\subsection*{B --- Analysis and candidate update}
B — Run its quality analysis and preserve the result.

\subsection*{C --- Independent validation}
C — Use the result as a prerequisite or independent validation for downstream fitted stages.

\subsection*{Measured / fitted values}
\begin{itemize}\setlength\itemsep{0pt}
\item \texttt{No scalar result has been written yet.}
\end{itemize}

\begin{center}
\includegraphics[width=0.78\textwidth,height=0.32\textheight,keepaspectratio]{../plots/stages/steering\_command\_audit.png}\\
\textit{Stage plot: Steering command-chain audit}
\end{center}

\vfill
\textbf{Raw data:} \texttt{../00\_command\_chain\_audit/}\\
\textbf{Stage report:} \texttt{stage\_reports/steering\_command\_audit.md}
\clearpage
\small
\section*{Direct physical vehicle metrology}
\addcontentsline{toc}{section}{Direct physical vehicle metrology}
\textbf{Stage key:} \texttt{physical\_metrology}\\
\textbf{Status:} pending\\

\subsection*{A --- Capture}
A — Measure the race-ready car using the session-local physical-measurements sheet.

\subsection*{B --- Analysis and candidate update}
B — Check SI units and mass/axle-load closure; derive l\_f/l\_r, freeze LiDAR/rear-axle/IMU reference geometry, apply the reversible dynamic-model geometry patch, and rebuild.

\subsection*{C --- Independent validation}
C — This is direct metrology, not a fitted time-series model; later independent motion stages use the rebuilt geometry and it remains an audited prerequisite for lateral identification.

\subsection*{Measured / fitted values}
\begin{itemize}\setlength\itemsep{0pt}
\item \texttt{No scalar result has been written yet.}
\end{itemize}

\textit{No plot has been generated for this stage yet.}

\vfill
\textbf{Raw data:} \texttt{../00a\_physical\_metrology/}\\
\textbf{Stage report:} \texttt{stage\_reports/physical\_metrology.md}
\clearpage
\small
\section*{LiDAR/IMU observability check}
\addcontentsline{toc}{section}{LiDAR/IMU observability check}
\textbf{Stage key:} \texttt{steering\_observability}\\
\textbf{Status:} pending\\

\subsection*{A --- Capture}
A — Record the stage-specific data.

\subsection*{B --- Analysis and candidate update}
B — Run its quality analysis and preserve the result.

\subsection*{C --- Independent validation}
C — Use the result as a prerequisite or independent validation for downstream fitted stages.

\subsection*{Measured / fitted values}
\begin{itemize}\setlength\itemsep{0pt}
\item \texttt{No scalar result has been written yet.}
\end{itemize}

\begin{center}
\includegraphics[width=0.78\textwidth,height=0.32\textheight,keepaspectratio]{../plots/stages/steering\_observability.png}\\
\textit{Stage plot: LiDAR/IMU observability check}
\end{center}

\vfill
\textbf{Raw data:} \texttt{../03\_sensor\_observability/}\\
\textbf{Stage report:} \texttt{stage\_reports/steering\_observability.md}
\clearpage
\small
\section*{Steering offset: training capture and fit}
\addcontentsline{toc}{section}{Steering offset: training capture and fit}
\textbf{Stage key:} \texttt{steering\_centre}\\
\textbf{Status:} pending\\

\subsection*{A --- Capture}
A — Record independent training captures.

\subsection*{B --- Analysis and candidate update}
B — Fit, gate, apply the candidate patch, and rebuild inside the reversible transaction.

\subsection*{C --- Independent validation}
C — The following hold-out/validation stage evaluates the applied candidate on new data.

\subsection*{Measured / fitted values}
\begin{itemize}\setlength\itemsep{0pt}
\item \texttt{No scalar result has been written yet.}
\end{itemize}

\begin{center}
\includegraphics[width=0.78\textwidth,height=0.32\textheight,keepaspectratio]{../plots/stages/steering\_centre.png}\\
\textit{Stage plot: Steering offset: training capture and fit}
\end{center}

\vfill
\textbf{Raw data:} \texttt{../01\_zero\_curvature\_centre/}\\
\textbf{Stage report:} \texttt{stage\_reports/steering\_centre.md}
\clearpage
\small
\section*{Steering offset: independent straightness validation}
\addcontentsline{toc}{section}{Steering offset: independent straightness validation}
\textbf{Stage key:} \texttt{steering\_centre\_validation}\\
\textbf{Status:} pending\\

\subsection*{A --- Capture}
A — Record a fresh validation data set after the preceding candidate was applied.

\subsection*{B --- Analysis and candidate update}
B — Analyse against the already fitted candidate; do not refit from the hold-out.

\subsection*{C --- Independent validation}
C — Pass only when the new-data gates and operator acceptance agree.

\subsection*{Measured / fitted values}
\begin{itemize}\setlength\itemsep{0pt}
\item \texttt{No scalar result has been written yet.}
\end{itemize}

\textit{No plot has been generated for this stage yet.}

\vfill
\textbf{Raw data:} \texttt{../01a\_zero\_curvature\_validation/}\\
\textbf{Stage report:} \texttt{stage\_reports/steering\_centre\_validation.md}
\clearpage
\small
\section*{Human steering-limit survey}
\addcontentsline{toc}{section}{Human steering-limit survey}
\textbf{Stage key:} \texttt{steering\_endstops}\\
\textbf{Status:} pending\\

\subsection*{A --- Capture}
A — Human records the last mechanically free steering limit on each side.

\subsection*{B --- Analysis and candidate update}
B — The safe raw-servo envelope is stored and checked against the validated centre.

\subsection*{C --- Independent validation}
C — This is a manual safety prerequisite, not a statistical fitted model.

\subsection*{Measured / fitted values}
\begin{itemize}\setlength\itemsep{0pt}
\item \texttt{No scalar result has been written yet.}
\end{itemize}

\begin{center}
\includegraphics[width=0.78\textwidth,height=0.32\textheight,keepaspectratio]{../plots/stages/steering\_endstops.png}\\
\textit{Stage plot: Human steering-limit survey}
\end{center}

\vfill
\textbf{Raw data:} \texttt{../02\_physical\_endstops/}\\
\textbf{Stage report:} \texttt{stage\_reports/steering\_endstops.md}
\clearpage
\small
\section*{Servo-to-effective-steering training}
\addcontentsline{toc}{section}{Servo-to-effective-steering training}
\textbf{Stage key:} \texttt{steering\_static\_training}\\
\textbf{Status:} pending\\

\subsection*{A --- Capture}
A — Record independent training captures.

\subsection*{B --- Analysis and candidate update}
B — Fit, gate, apply the candidate patch, and rebuild inside the reversible transaction.

\subsection*{C --- Independent validation}
C — The following hold-out/validation stage evaluates the applied candidate on new data.

\subsection*{Measured / fitted values}
\begin{itemize}\setlength\itemsep{0pt}
\item \texttt{No scalar result has been written yet.}
\end{itemize}

\begin{center}
\includegraphics[width=0.78\textwidth,height=0.32\textheight,keepaspectratio]{../plots/stages/steering\_static\_training.png}\\
\textit{Stage plot: Servo-to-effective-steering training}
\end{center}

\vfill
\textbf{Raw data:} \texttt{../04\_static\_map\_training/}\\
\textbf{Stage report:} \texttt{stage\_reports/steering\_static\_training.md}
\clearpage
\small
\section*{Servo-to-effective-steering hold-out validation}
\addcontentsline{toc}{section}{Servo-to-effective-steering hold-out validation}
\textbf{Stage key:} \texttt{steering\_static\_holdout}\\
\textbf{Status:} pending\\

\subsection*{A --- Capture}
A — Record a fresh validation data set after the preceding candidate was applied.

\subsection*{B --- Analysis and candidate update}
B — Analyse against the already fitted candidate; do not refit from the hold-out.

\subsection*{C --- Independent validation}
C — Pass only when the new-data gates and operator acceptance agree.

\subsection*{Measured / fitted values}
\begin{itemize}\setlength\itemsep{0pt}
\item \texttt{No scalar result has been written yet.}
\end{itemize}

\begin{center}
\includegraphics[width=0.78\textwidth,height=0.32\textheight,keepaspectratio]{../plots/stages/steering\_static\_holdout.png}\\
\textit{Stage plot: Servo-to-effective-steering hold-out validation}
\end{center}

\vfill
\textbf{Raw data:} \texttt{../05\_static\_map\_holdout/}\\
\textbf{Stage report:} \texttt{stage\_reports/steering\_static\_holdout.md}
\clearpage
\small
\section*{Steering response: training characterisation}
\addcontentsline{toc}{section}{Steering response: training characterisation}
\textbf{Stage key:} \texttt{steering\_response}\\
\textbf{Status:} pending\\

\subsection*{A --- Capture}
A — Record the stage-specific candidate/training command grid.

\subsection*{B --- Analysis and candidate update}
B — Estimate timing/routing metrics and preserve the candidate evidence; no hidden refit is made on later C data.

\subsection*{C --- Independent validation}
C — The following explicitly named hold-out stage records different commands and evaluates the frozen result.

\subsection*{Measured / fitted values}
\begin{itemize}\setlength\itemsep{0pt}
\item \texttt{No scalar result has been written yet.}
\end{itemize}

\begin{center}
\includegraphics[width=0.78\textwidth,height=0.32\textheight,keepaspectratio]{../plots/stages/steering\_response.png}\\
\textit{Stage plot: Steering response: training characterisation}
\end{center}

\vfill
\textbf{Raw data:} \texttt{../06\_command\_to\_curvature\_response/}\\
\textbf{Stage report:} \texttt{stage\_reports/steering\_response.md}
\clearpage
\small
\section*{Steering response: independent hold-out}
\addcontentsline{toc}{section}{Steering response: independent hold-out}
\textbf{Stage key:} \texttt{steering\_response\_validation}\\
\textbf{Status:} pending\\

\subsection*{A --- Capture}
A — Record a fresh validation data set after the preceding candidate was applied.

\subsection*{B --- Analysis and candidate update}
B — Analyse against the already fitted candidate; do not refit from the hold-out.

\subsection*{C --- Independent validation}
C — Pass only when the new-data gates and operator acceptance agree.

\subsection*{Measured / fitted values}
\begin{itemize}\setlength\itemsep{0pt}
\item \texttt{No scalar result has been written yet.}
\end{itemize}

\textit{No plot has been generated for this stage yet.}

\vfill
\textbf{Raw data:} \texttt{../06a\_command\_to\_curvature\_response\_validation/}\\
\textbf{Stage report:} \texttt{stage\_reports/steering\_response\_validation.md}
\clearpage
\small
\section*{Motor command-chain audit}
\addcontentsline{toc}{section}{Motor command-chain audit}
\textbf{Stage key:} \texttt{motor\_command\_audit}\\
\textbf{Status:} pending\\

\subsection*{A --- Capture}
A — Record the stage-specific data.

\subsection*{B --- Analysis and candidate update}
B — Run its quality analysis and preserve the result.

\subsection*{C --- Independent validation}
C — Use the result as a prerequisite or independent validation for downstream fitted stages.

\subsection*{Measured / fitted values}
\begin{itemize}\setlength\itemsep{0pt}
\item \texttt{No scalar result has been written yet.}
\end{itemize}

\textit{No plot has been generated for this stage yet.}

\vfill
\textbf{Raw data:} \texttt{../00b\_motor\_command\_chain\_audit/}\\
\textbf{Stage report:} \texttt{stage\_reports/motor\_command\_audit.md}
\clearpage
\small
\section*{Longitudinal sensor observability}
\addcontentsline{toc}{section}{Longitudinal sensor observability}
\textbf{Stage key:} \texttt{longitudinal\_observability}\\
\textbf{Status:} pending\\

\subsection*{A --- Capture}
A — Capture a stationary diagnostic and repeated low/medium/high-speed straight passes.

\subsection*{B --- Analysis and candidate update}
B — Gate the robust moving LiDAR-window measurement quality and write the stationary noise diagnostic.

\subsection*{C --- Independent validation}
C — Later ERPM training is blocked unless this independent moving-sensor preflight passed.

\subsection*{Measured / fitted values}
\begin{itemize}\setlength\itemsep{0pt}
\item \texttt{No scalar result has been written yet.}
\end{itemize}

\textit{No plot has been generated for this stage yet.}

\vfill
\textbf{Raw data:} \texttt{../01\_longitudinal\_observability/}\\
\textbf{Stage report:} \texttt{stage\_reports/longitudinal\_observability.md}
\clearpage
\small
\section*{Low-speed launch characterisation}
\addcontentsline{toc}{section}{Low-speed launch characterisation}
\textbf{Stage key:} \texttt{low\_speed\_launch}\\
\textbf{Status:} pending\\

\subsection*{A --- Capture}
A — Capture repeated raw-ERPM launches at the low-speed grid.

\subsection*{B --- Analysis and candidate update}
B — Report condition coverage and the first repeatable LiDAR ground-motion threshold; do not prematurely fit the full speed map.

\subsection*{C --- Independent validation}
C — The later rebuilt VEL\_TO\_ERPM pipeline audit checks the resulting slow-start setting on new commands.

\subsection*{Measured / fitted values}
\begin{itemize}\setlength\itemsep{0pt}
\item \texttt{No scalar result has been written yet.}
\end{itemize}

\textit{No plot has been generated for this stage yet.}

\vfill
\textbf{Raw data:} \texttt{../02\_low\_speed\_launch/}\\
\textbf{Stage report:} \texttt{stage\_reports/low\_speed\_launch.md}
\clearpage
\small
\section*{ERPM-to-ground-speed training}
\addcontentsline{toc}{section}{ERPM-to-ground-speed training}
\textbf{Stage key:} \texttt{erpm\_map\_training}\\
\textbf{Status:} pending\\

\subsection*{A --- Capture}
A — Record independent training captures.

\subsection*{B --- Analysis and candidate update}
B — Fit, gate, apply the candidate patch, and rebuild inside the reversible transaction.

\subsection*{C --- Independent validation}
C — The following hold-out/validation stage evaluates the applied candidate on new data.

\subsection*{Measured / fitted values}
\begin{itemize}\setlength\itemsep{0pt}
\item \texttt{No scalar result has been written yet.}
\end{itemize}

\textit{No plot has been generated for this stage yet.}

\vfill
\textbf{Raw data:} \texttt{../03\_raw\_erpm\_map\_training/}\\
\textbf{Stage report:} \texttt{stage\_reports/erpm\_map\_training.md}
\clearpage
\small
\section*{ERPM-to-ground-speed hold-out validation}
\addcontentsline{toc}{section}{ERPM-to-ground-speed hold-out validation}
\textbf{Stage key:} \texttt{erpm\_map\_holdout}\\
\textbf{Status:} pending\\

\subsection*{A --- Capture}
A — Record a fresh validation data set after the preceding candidate was applied.

\subsection*{B --- Analysis and candidate update}
B — Analyse against the already fitted candidate; do not refit from the hold-out.

\subsection*{C --- Independent validation}
C — Pass only when the new-data gates and operator acceptance agree.

\subsection*{Measured / fitted values}
\begin{itemize}\setlength\itemsep{0pt}
\item \texttt{No scalar result has been written yet.}
\end{itemize}

\textit{No plot has been generated for this stage yet.}

\vfill
\textbf{Raw data:} \texttt{../04\_raw\_erpm\_map\_holdout/}\\
\textbf{Stage report:} \texttt{stage\_reports/erpm\_map\_holdout.md}
\clearpage
\small
\section*{Velocity-to-ERPM command-path validation}
\addcontentsline{toc}{section}{Velocity-to-ERPM command-path validation}
\textbf{Stage key:} \texttt{vel\_to\_erpm\_audit}\\
\textbf{Status:} pending\\

\subsection*{A --- Capture}
A — Record the stage-specific data.

\subsection*{B --- Analysis and candidate update}
B — Run its quality analysis and preserve the result.

\subsection*{C --- Independent validation}
C — Use the result as a prerequisite or independent validation for downstream fitted stages.

\subsection*{Measured / fitted values}
\begin{itemize}\setlength\itemsep{0pt}
\item \texttt{No scalar result has been written yet.}
\end{itemize}

\textit{No plot has been generated for this stage yet.}

\vfill
\textbf{Raw data:} \texttt{../05\_vel\_to\_erpm\_pipeline\_audit/}\\
\textbf{Stage report:} \texttt{stage\_reports/vel\_to\_erpm\_audit.md}
\clearpage
\small
\section*{ERPM response characterisation}
\addcontentsline{toc}{section}{ERPM response characterisation}
\textbf{Stage key:} \texttt{erpm\_response}\\
\textbf{Status:} pending\\

\subsection*{A --- Capture}
A — Record the stage-specific candidate/training command grid.

\subsection*{B --- Analysis and candidate update}
B — Estimate timing/routing metrics and preserve the candidate evidence; no hidden refit is made on later C data.

\subsection*{C --- Independent validation}
C — The following explicitly named hold-out stage records different commands and evaluates the frozen result.

\subsection*{Measured / fitted values}
\begin{itemize}\setlength\itemsep{0pt}
\item \texttt{No scalar result has been written yet.}
\end{itemize}

\textit{No plot has been generated for this stage yet.}

\vfill
\textbf{Raw data:} \texttt{../06\_raw\_erpm\_response/}\\
\textbf{Stage report:} \texttt{stage\_reports/erpm\_response.md}
\clearpage
\small
\section*{ERPM response independent hold-out}
\addcontentsline{toc}{section}{ERPM response independent hold-out}
\textbf{Stage key:} \texttt{erpm\_response\_validation}\\
\textbf{Status:} pending\\

\subsection*{A --- Capture}
A — Record a fresh validation data set after the preceding candidate was applied.

\subsection*{B --- Analysis and candidate update}
B — Analyse against the already fitted candidate; do not refit from the hold-out.

\subsection*{C --- Independent validation}
C — Pass only when the new-data gates and operator acceptance agree.

\subsection*{Measured / fitted values}
\begin{itemize}\setlength\itemsep{0pt}
\item \texttt{No scalar result has been written yet.}
\end{itemize}

\textit{No plot has been generated for this stage yet.}

\vfill
\textbf{Raw data:} \texttt{../06a\_raw\_erpm\_response\_validation/}\\
\textbf{Stage report:} \texttt{stage\_reports/erpm\_response\_validation.md}
\clearpage
\small
\section*{Coast-down drag fit}
\addcontentsline{toc}{section}{Coast-down drag fit}
\textbf{Stage key:} \texttt{coastdown}\\
\textbf{Status:} pending\\

\subsection*{A --- Capture}
A — Record independent training captures.

\subsection*{B --- Analysis and candidate update}
B — Fit, gate, apply the candidate patch, and rebuild inside the reversible transaction.

\subsection*{C --- Independent validation}
C — The following hold-out/validation stage evaluates the applied candidate on new data.

\subsection*{Measured / fitted values}
\begin{itemize}\setlength\itemsep{0pt}
\item \texttt{No scalar result has been written yet.}
\end{itemize}

\textit{No plot has been generated for this stage yet.}

\vfill
\textbf{Raw data:} \texttt{../07\_coastdown/}\\
\textbf{Stage report:} \texttt{stage\_reports/coastdown.md}
\clearpage
\small
\section*{Coast-down drag independent hold-out}
\addcontentsline{toc}{section}{Coast-down drag independent hold-out}
\textbf{Stage key:} \texttt{coastdown\_validation}\\
\textbf{Status:} pending\\

\subsection*{A --- Capture}
A — Record a fresh validation data set after the preceding candidate was applied.

\subsection*{B --- Analysis and candidate update}
B — Analyse against the already fitted candidate; do not refit from the hold-out.

\subsection*{C --- Independent validation}
C — Pass only when the new-data gates and operator acceptance agree.

\subsection*{Measured / fitted values}
\begin{itemize}\setlength\itemsep{0pt}
\item \texttt{No scalar result has been written yet.}
\end{itemize}

\textit{No plot has been generated for this stage yet.}

\vfill
\textbf{Raw data:} \texttt{../07a\_coastdown\_validation/}\\
\textbf{Stage report:} \texttt{stage\_reports/coastdown\_validation.md}
\clearpage
\small
\section*{Current-to-acceleration training}
\addcontentsline{toc}{section}{Current-to-acceleration training}
\textbf{Stage key:} \texttt{current\_training}\\
\textbf{Status:} pending\\

\subsection*{A --- Capture}
A — Record independent training captures.

\subsection*{B --- Analysis and candidate update}
B — Fit, gate, apply the candidate patch, and rebuild inside the reversible transaction.

\subsection*{C --- Independent validation}
C — The following hold-out/validation stage evaluates the applied candidate on new data.

\subsection*{Measured / fitted values}
\begin{itemize}\setlength\itemsep{0pt}
\item \texttt{No scalar result has been written yet.}
\end{itemize}

\textit{No plot has been generated for this stage yet.}

\vfill
\textbf{Raw data:} \texttt{../08\_raw\_current\_training/}\\
\textbf{Stage report:} \texttt{stage\_reports/current\_training.md}
\clearpage
\small
\section*{Current-to-acceleration hold-out validation}
\addcontentsline{toc}{section}{Current-to-acceleration hold-out validation}
\textbf{Stage key:} \texttt{current\_holdout}\\
\textbf{Status:} pending\\

\subsection*{A --- Capture}
A — Record a fresh validation data set after the preceding candidate was applied.

\subsection*{B --- Analysis and candidate update}
B — Analyse against the already fitted candidate; do not refit from the hold-out.

\subsection*{C --- Independent validation}
C — Pass only when the new-data gates and operator acceptance agree.

\subsection*{Measured / fitted values}
\begin{itemize}\setlength\itemsep{0pt}
\item \texttt{No scalar result has been written yet.}
\end{itemize}

\textit{No plot has been generated for this stage yet.}

\vfill
\textbf{Raw data:} \texttt{../09\_raw\_current\_holdout/}\\
\textbf{Stage report:} \texttt{stage\_reports/current\_holdout.md}
\clearpage
\small
\section*{Acceleration-command interface training audit}
\addcontentsline{toc}{section}{Acceleration-command interface training audit}
\textbf{Stage key:} \texttt{accel\_interface}\\
\textbf{Status:} pending\\

\subsection*{A --- Capture}
A — Record the stage-specific candidate/training command grid.

\subsection*{B --- Analysis and candidate update}
B — Estimate timing/routing metrics and preserve the candidate evidence; no hidden refit is made on later C data.

\subsection*{C --- Independent validation}
C — The following explicitly named hold-out stage records different commands and evaluates the frozen result.

\subsection*{Measured / fitted values}
\begin{itemize}\setlength\itemsep{0pt}
\item \texttt{No scalar result has been written yet.}
\end{itemize}

\textit{No plot has been generated for this stage yet.}

\vfill
\textbf{Raw data:} \texttt{../10\_accel\_to\_current\_interface/}\\
\textbf{Stage report:} \texttt{stage\_reports/accel\_interface.md}
\clearpage
\small
\section*{Acceleration-command interface independent hold-out}
\addcontentsline{toc}{section}{Acceleration-command interface independent hold-out}
\textbf{Stage key:} \texttt{accel\_interface\_validation}\\
\textbf{Status:} pending\\

\subsection*{A --- Capture}
A — Record a fresh validation data set after the preceding candidate was applied.

\subsection*{B --- Analysis and candidate update}
B — Analyse against the already fitted candidate; do not refit from the hold-out.

\subsection*{C --- Independent validation}
C — Pass only when the new-data gates and operator acceptance agree.

\subsection*{Measured / fitted values}
\begin{itemize}\setlength\itemsep{0pt}
\item \texttt{No scalar result has been written yet.}
\end{itemize}

\textit{No plot has been generated for this stage yet.}

\vfill
\textbf{Raw data:} \texttt{../10a\_accel\_to\_current\_interface\_validation/}\\
\textbf{Stage report:} \texttt{stage\_reports/accel\_interface\_validation.md}
\clearpage
\small
\section*{Selected odometry model: fresh steady-speed validation}
\addcontentsline{toc}{section}{Selected odometry model: fresh steady-speed validation}
\textbf{Stage key:} \texttt{odometry\_candidate\_velocity\_validation}\\
\textbf{Status:} pending\\

\subsection*{A --- Capture}
A — Record a fresh validation data set after the preceding candidate was applied.

\subsection*{B --- Analysis and candidate update}
B — Analyse against the already fitted candidate; do not refit from the hold-out.

\subsection*{C --- Independent validation}
C — Pass only when the new-data gates and operator acceptance agree.

\subsection*{Measured / fitted values}
\begin{itemize}\setlength\itemsep{0pt}
\item \texttt{No scalar result has been written yet.}
\end{itemize}

\textit{No plot has been generated for this stage yet.}

\vfill
\textbf{Raw data:} \texttt{../11\_candidate\_velocity\_verification/}\\
\textbf{Stage report:} \texttt{stage\_reports/odometry\_candidate\_velocity\_validation.md}
\clearpage
\small
\section*{Selected odometry model: fresh transient validation}
\addcontentsline{toc}{section}{Selected odometry model: fresh transient validation}
\textbf{Stage key:} \texttt{odometry\_candidate\_accel\_validation}\\
\textbf{Status:} pending\\

\subsection*{A --- Capture}
A — Record a fresh validation data set after the preceding candidate was applied.

\subsection*{B --- Analysis and candidate update}
B — Analyse against the already fitted candidate; do not refit from the hold-out.

\subsection*{C --- Independent validation}
C — Pass only when the new-data gates and operator acceptance agree.

\subsection*{Measured / fitted values}
\begin{itemize}\setlength\itemsep{0pt}
\item \texttt{No scalar result has been written yet.}
\end{itemize}

\textit{No plot has been generated for this stage yet.}

\vfill
\textbf{Raw data:} \texttt{../12\_candidate\_accel\_verification/}\\
\textbf{Stage report:} \texttt{stage\_reports/odometry\_candidate\_accel\_validation.md}
\clearpage
\small
\section*{Effective tyre stiffness: quasi-steady training arcs}
\addcontentsline{toc}{section}{Effective tyre stiffness: quasi-steady training arcs}
\textbf{Stage key:} \texttt{lateral\_stiffness\_training}\\
\textbf{Status:} pending\\

\subsection*{A --- Capture}
A — Record balanced, quasi-steady left/right speed-and-steering arcs after physical metrology.

\subsection*{B --- Analysis and candidate update}
B — Fit effective low-slip front/rear stiffness plus the dynamic-model steering scale, apply the reversible bicycle-model patch, and retain a diagnostic nonlinear candidate.

\subsection*{C --- Independent validation}
C — The following hold-out uses distinct arcs and blocks promotion if the frozen scale/stiffness model fails or a nonlinear tyre model is materially better.

\subsection*{Measured / fitted values}
\begin{itemize}\setlength\itemsep{0pt}
\item \texttt{No scalar result has been written yet.}
\end{itemize}

\textit{No plot has been generated for this stage yet.}

\vfill
\textbf{Raw data:} \texttt{../12\_quasi\_steady\_lateral\_training/}\\
\textbf{Stage report:} \texttt{stage\_reports/lateral\_stiffness\_training.md}
\clearpage
\small
\section*{Effective tyre stiffness: independent hold-out arcs}
\addcontentsline{toc}{section}{Effective tyre stiffness: independent hold-out arcs}
\textbf{Stage key:} \texttt{lateral\_stiffness\_validation}\\
\textbf{Status:} pending\\

\subsection*{A --- Capture}
A — Record a fresh validation data set after the preceding candidate was applied.

\subsection*{B --- Analysis and candidate update}
B — Analyse against the already fitted candidate; do not refit from the hold-out.

\subsection*{C --- Independent validation}
C — Pass only when the new-data gates and operator acceptance agree.

\subsection*{Measured / fitted values}
\begin{itemize}\setlength\itemsep{0pt}
\item \texttt{No scalar result has been written yet.}
\end{itemize}

\textit{No plot has been generated for this stage yet.}

\vfill
\textbf{Raw data:} \texttt{../12a\_quasi\_steady\_lateral\_validation/}\\
\textbf{Stage report:} \texttt{stage\_reports/lateral\_stiffness\_validation.md}
\clearpage
\end{document}
